\begin{center}
\mbox{}

\bigskip

\LARGE\textsc{Prier selon le tabernacle}

\vspace{2ex plus \fill}

\large{d'après le pasteur Jon Courson}

\vspace{3ex plus \fill}

\includegraphics[width=0.2\textwidth]{colombe.pdf}

\vspace{0ex plus \fill}


\mbox{}
\end{center}

\pagestyle{empty}
\setlength{\parindent}{0em}

\section{Les portes (remerciements)}

Entrer dans ses portes avec des actions de grâce (Psaume~100.4).

Remercier Dieu pour tous ses bienfaits, pour ce qu'Il donne ! C'est la \og répression de la dépression \fg{}!

 
\section{Le parvis (louange)}

Entrer dans ses parvis avec la louange (Psaume~100.4).

Le louer, le célébrer pour qui Il est et pour ce qu'Il a fait et ce qu'Il fait\dots{} C'est la répression de l'oppression! L'ennemi s'enfuit!

Tactique du roi Josaphat:  la louange en tête du combat!

\begin{quote}
\og Puis, il tint conseil avec le peuple et désigna des chantres pour l’Éternel, afin de faire entendre la louange avec des ornements sacrés, lorsqu’ils sortiraient devant l’armée. Ils disaient: Célébrez l’Éternel, car sa bienveillance dure à toujours! \fg{}

\raggedleft\small 2~Chroniques~20.21
\end{quote}

 
\section{Autel de bronze (confession)}

Sacrifices offerts pour mes péchés\dots

La confession de mes péchés a deux effets principaux:

\begin{itemize}
\item Libération de mes péchés: le péché perd son pouvoir!
\item Appréciation pour mon Sauveur.
\end{itemize}

 
\section{Lavoir (méditation de la parole)}

Se laver avec la Parole de Dieu.

\begin{quote}
Le Christ a aimé l’Église et s’est livré lui-même pour elle, afin de la sanctifier après l’avoir purifiée par l’eau et la parole.

\raggedleft\small Ephésiens~5.25--26
\end{quote}


 
\section{Table des pains (besoins\\personnels)}

Pains de Proposition.

Finances, choses physiques, santé, etc.

 

\section{Chandeliers d'or (ministères)}

Nous sommes la Lumière du Monde.

Faire briller l'Évangile.

Prier pour tous les ministères auxquels Dieu nous fait penser:

\begin{itemize}
  \item dans l'église;
  \item dans le pays;
  \item à l'international;
  \item missions;
  \item etc.
\end{itemize}
 

\section{Autel des parfums (intercession)}

Faire monter des intercessions comme un parfum de bonne odeur.

Prier pour les besoins des autres autour de nous.

 

\section{Saint des Saints (adoration)}

Adorer Dieu.
